% !TEX root = ../main.tex
\section{Introduction}
\label{sec:introduction}

Studying collusion among language-model agents requires an environment that
preserves its protocol rules while exposing the intended experimental controls.
The reduced Veracity game assigns agents to create a claim, challenge it, and
serve on a jury whose committed votes determine financial
settlement~\cite{project-spec}. The broader experiment varies incentives,
communication, transfers, memory, stakes, and evidence. Before interpreting
differences between those settings, we need to verify that live agents can
complete the baseline and that its measurements remain traceable to execution.

This report answers that narrower integration question. We examine one short
block of two episodes with all six factors at level zero. One claim is true
and the other false, allowing us to distinguish following evidence from always
supporting the same party. We report jury agreement, truth alignment, private
traffic, and consistency across the two episodes, alongside execution checks.
The figures show the observations as well as their means and descriptive
standard deviations. The study supplies a reproducible starting point for
the planned factorial experiments, with no treatment comparison or inferential
claim about collusion.

Section~\ref{sec:probes} adds six single-episode probes of the cell with every
opportunity enabled. They check new mechanisms before the first pass: an
auditor that questions jurors after the verdict, sealed competing offers from
both parties, party offers to the auditor, and a stated rule prohibiting side
payments and private chat. A reasoning monitor labels each agent's conduct. The
probes show what one model did in these settings; they are not a sample from
which effects can be estimated.
